% GENERATED FILE. Edit canonical lessons, then run npm run generate:books.
% canonical-source-hash: fnv1a64:86a4ee1b264a3c19
% canonical-lessons: HI-C422-jhopri, HI-C422-kothi, HI-C422-makan, HI-C422-ahata, HI-C422-chappar

\chapter{Houses Large and Small}
\label{ch:hi-houses-large-and-small}
\hlchaptermodality{422}

\begin{chapteropening}
\textbf{By the end of this chapter:} \emph{I can describe a hut, a bungalow, a house, a compound and a thatched roof.}

Complete the last lesson of chapter 422: I can describe a hut, a bungalow, a house, a compound and a thatched roof.
\end{chapteropening}

\section[\texorpdfstring{\hi{झोंपड़ी} (jhõpṛī)}{झोंपड़ी (jhõpṛī)}]{\hi{झोंपड़ी} (jhõpṛī) — a hut}

\label{lesson:HI-C422-jhopri}



\begin{quote}
Before the new one: say the Hindi for next door, then the Hindi for the countryside.
\end{quote}

\subsection*{You'll want to know: \hi{झोंपड़ी}}
\textbf{\hi{झोंपड़ी}} — \emph{jhõpṛī} — \textquotedblleft{}a hut\textquotedblright{}.

\begin{quote}
\hi{खेत} \hi{के} \hi{पास} \hi{एक} \hi{झोंपड़ी} \hi{है।}
\end{quote}

\begin{quote}
\emph{khet ke pās ek jhõpṛī hai.}
\end{quote}

\begin{quote}
There is a hut near the field.
\end{quote}

Feminine.

\begin{cousinweb}
\textbf{\hi{झोंपड़ी}} (\emph{jhõpṛī}) goes back to \emph{*jhuppa}, \textquotedblleft{}a hut\textquotedblright{}, a form CDIAL reconstructs (CDIAL 5403), which lists Hindi \emph{jhõpṛī}.
\end{cousinweb}

\begin{grammarlens}[title={What you've built}]
One more word for getting around.
\end{grammarlens}

\subsection*{Your turn}
\begin{itemize}
\raggedright
  \item \emph{Say it:} \emph{jhõpṛī}
  \item \emph{Say it:} \emph{jhõpṛī}, once more
  \item \emph{Say it:} \emph{khet ke pās ek jhõpṛī hai.} — There is a hut near the field.
  \item \emph{Recall:} say the Hindi for next door, then the Hindi for the countryside, then say \emph{jhõpṛī} again
\end{itemize}

\begin{tcolorbox}[breakable,colback=teal!4,colframe=teal!35!black,title={Before you move on}]
What does \hi{झोंपड़ी} mean? (\textquotedblleft{}A hut\textquotedblright{}.) Say it once more.
\end{tcolorbox}
\section[\texorpdfstring{\hi{कोठी} (koṭhī)}{कोठी (koṭhī)}]{\hi{कोठी} (koṭhī) — a large house, a bungalow}

\label{lesson:HI-C422-kothi}



\begin{quote}
Before the new one: say the Hindi for the countryside, then the Hindi for a hut.
\end{quote}

\subsection*{You'll want to know: \hi{कोठी}}
\textbf{\hi{कोठी}} — \emph{koṭhī} — \textquotedblleft{}a large house, a bungalow\textquotedblright{}.

\begin{quote}
\hi{उस} \hi{बड़ी} \hi{कोठी} \hi{में} \hi{एक} \hi{अमीर} \hi{परिवार} \hi{रहता} \hi{है।}
\end{quote}

\begin{quote}
\emph{us baṛī koṭhī meṁ ek amīr parivār rahtā hai.}
\end{quote}

\begin{quote}
A rich family lives in that big bungalow.
\end{quote}

Feminine.

\begin{cousinweb}
\textbf{\hi{कोठी}} (\emph{koṭhī}) is a tadbhava of Sanskrit \emph{koṣṭha}, \textquotedblleft{}a granary, a storeroom\textquotedblright{} (CDIAL 3546), which lists Hindi \emph{koṭhī}, \textquotedblleft{}granary, large house\textquotedblright{}: the strong room grew into the whole mansion.
\end{cousinweb}

\begin{grammarlens}[title={What you've built}]
One more word for getting around.
\end{grammarlens}

\subsection*{Your turn}
\begin{itemize}
\raggedright
  \item \emph{Say it:} \emph{koṭhī}
  \item \emph{Say it:} \emph{koṭhī}, once more
  \item \emph{Say it:} \emph{us baṛī koṭhī meṁ ek amīr parivār rahtā hai.} — A rich family lives in that big bungalow.
  \item \emph{Recall:} say the Hindi for the countryside, then the Hindi for a hut, then say \emph{koṭhī} again
\end{itemize}

\begin{tcolorbox}[breakable,colback=teal!4,colframe=teal!35!black,title={Before you move on}]
What does \hi{कोठी} mean? (\textquotedblleft{}A large house, a bungalow\textquotedblright{}.) Say it once more.
\end{tcolorbox}
\section[\texorpdfstring{\hi{मकान} (makān)}{मकान (makān)}]{\hi{मकान} (makān) — a house (the building)}

\label{lesson:HI-C422-makan}



\begin{quote}
Before the new one: say the Hindi for a hut, then the Hindi for a bungalow.
\end{quote}

\subsection*{You'll want to know: \hi{मकान}}
\textbf{\hi{मकान}} — \emph{makān} — \textquotedblleft{}a house (the building)\textquotedblright{}.

\begin{quote}
\hi{हमने} \hi{किराये} \hi{पर} \hi{एक} \hi{मकान} \hi{लिया।}
\end{quote}

\begin{quote}
\emph{hamne kirāye par ek makān liyā.}
\end{quote}

\begin{quote}
We rented a house.
\end{quote}

Masculine.

\hi{घर} is home; \hi{मकान} is the building, the thing you rent, buy or build.

\begin{cousinweb}
\textbf{\hi{मकान}} (\emph{makān}) is usually traced to Arabic \emph{makān}, \textquotedblleft{}a place\textquotedblright{}, which came through Persian and narrowed in Hindi to a house as a building.
\end{cousinweb}

\begin{grammarlens}[title={What you've built}]
One more word for getting around.
\end{grammarlens}

\subsection*{Your turn}
\begin{itemize}
\raggedright
  \item \emph{Say it:} \emph{makān}
  \item \emph{Say it:} \emph{makān}, once more
  \item \emph{Say it:} \emph{hamne kirāye par ek makān liyā.} — We rented a house.
  \item \emph{Recall:} say the Hindi for a hut, then the Hindi for a bungalow, then say \emph{makān} again
\end{itemize}

\begin{tcolorbox}[breakable,colback=teal!4,colframe=teal!35!black,title={Before you move on}]
What does \hi{मकान} mean? (\textquotedblleft{}A house (the building)\textquotedblright{}.) Say it once more.
\end{tcolorbox}
\section[\texorpdfstring{\hi{अहाता} (ahātā)}{अहाता (ahātā)}]{\hi{अहाता} (ahātā) — a compound, an enclosure}

\label{lesson:HI-C422-ahata}



\begin{quote}
Before the new one: say the Hindi for a bungalow, then the Hindi for a house (the building).
\end{quote}

\subsection*{You'll want to know: \hi{अहाता}}
\textbf{\hi{अहाता}} — \emph{ahātā} — \textquotedblleft{}a compound, an enclosure\textquotedblright{}.

\begin{quote}
\hi{स्कूल} \hi{के} \hi{अहाते} \hi{में} \hi{बच्चे} \hi{खेल} \hi{रहे} \hi{थे।}
\end{quote}

\begin{quote}
\emph{skūl ke ahāte meṁ bacce khel rahe the.}
\end{quote}

\begin{quote}
The children were playing in the school compound.
\end{quote}

Masculine.

\begin{cousinweb}
\textbf{\hi{अहाता}} (\emph{ahātā}) is usually traced to Arabic \emph{iḥāṭa}, \textquotedblleft{}an enclosing, an enclosure\textquotedblright{}, through Persian: the walled ground round a house, a school or an office.
\end{cousinweb}

\begin{grammarlens}[title={What you've built}]
One more word for getting around.
\end{grammarlens}

\subsection*{Your turn}
\begin{itemize}
\raggedright
  \item \emph{Say it:} \emph{ahātā}
  \item \emph{Say it:} \emph{ahātā}, once more
  \item \emph{Say it:} \emph{skūl ke ahāte meṁ bacce khel rahe the.} — The children were playing in the school compound.
  \item \emph{Recall:} say the Hindi for a bungalow, then the Hindi for a house (the building), then say \emph{ahātā} again
\end{itemize}

\begin{tcolorbox}[breakable,colback=teal!4,colframe=teal!35!black,title={Before you move on}]
What does \hi{अहाता} mean? (\textquotedblleft{}A compound, an enclosure\textquotedblright{}.) Say it once more.
\end{tcolorbox}
\section[\texorpdfstring{\hi{छप्पर} (chappar)}{छप्पर (chappar)}]{\hi{छप्पर} (chappar) — a thatched roof}

\label{lesson:HI-C422-chappar}



\begin{quote}
Before the new one: say the Hindi for a house (the building), then the Hindi for a compound.
\end{quote}

\subsection*{You'll want to know: \hi{छप्पर}}
\textbf{\hi{छप्पर}} — \emph{chappar} — \textquotedblleft{}a thatched roof\textquotedblright{}.

\begin{quote}
\hi{बारिश} \hi{में} \hi{झोंपड़ी} \hi{का} \hi{छप्पर} \hi{टपकता} \hi{है।}
\end{quote}

\begin{quote}
\emph{bāriś meṁ jhõpṛī kā chappar ṭapaktā hai.}
\end{quote}

\begin{quote}
The hut's thatched roof leaks in the rain.
\end{quote}

Masculine.

\begin{cousinweb}
\textbf{\hi{छप्पर}} (\emph{chappar}) is a tadbhava of Sanskrit \emph{chattvara}, \textquotedblleft{}a house, a bower\textquotedblright{} (CDIAL 4976), which lists Hindi \emph{chappar}, \textquotedblleft{}a sloping thatch, a hut\textquotedblright{}.
\end{cousinweb}

\begin{grammarlens}[title={What you've built}]
One more word for getting around.
\end{grammarlens}

\subsection*{Your turn}
\begin{itemize}
\raggedright
  \item \emph{Say it:} \emph{chappar}
  \item \emph{Say it:} \emph{chappar}, once more
  \item \emph{Say it:} \emph{bāriś meṁ jhõpṛī kā chappar ṭapaktā hai.} — The hut's thatched roof leaks in the rain.
  \item \emph{Recall:} say the Hindi for a house (the building), then the Hindi for a compound, then say \emph{chappar} again
\end{itemize}

\begin{tcolorbox}[breakable,colback=teal!4,colframe=teal!35!black,title={Before you move on}]
What does \hi{छप्पर} mean? (\textquotedblleft{}A thatched roof\textquotedblright{}.) Say it once more.
\end{tcolorbox}
